% Source problem: https://www.ap.t.u-tokyo.ac.jp/themes/ap_t_u-tokyo/graduate/pdf/2021ap_exam.pdf
\documentclass[11pt,a4paper]{article}
\usepackage[margin=24mm]{geometry}
\usepackage{fontspec}
\usepackage{xeCJK}
\setCJKmainfont{HaranoAjiMincho-Regular.otf}
\setCJKsansfont{HaranoAjiGothic-Medium.otf}
\usepackage{amsmath,amssymb}
\setlength{\parindent}{1em}
\setlength{\parskip}{0.35em}
\title{物理工学専攻 院試解答集\\2021年度 物理学 第3問}
\author{政岡凜太郎}
\date{}
\begin{document}
\maketitle

\section*{[1.1] 箱中の運動量状態数}
周期境界条件では各運動量成分が $p_j=2\pi\hbar n_j/L$ と量子化される。
運動量空間の1状態当たりの体積は $(2\pi\hbar/L)^3$ だから、単位体積当たり
\begin{equation}
 \boxed{g_p=\left(\frac{L}{2\pi\hbar}\right)^3}
\end{equation}
個の状態がある。粒子のスピン自由度は数えない。

\section*{[1.2] 三次元箱の状態密度}
運動量の大きさが $p$ から $p+dp$ の殻の体積は $4\pi p^2dp$ であり、
$\epsilon=p^2/(2m)$ より $dp/d\epsilon=m/p$ である。従って
\begin{align}
 D_3(\epsilon)d\epsilon
 &=4\pi p^2\frac{dp}{d\epsilon}d\epsilon
   \left(\frac{L}{2\pi\hbar}\right)^3,\\
 D_3(\epsilon)
 &=\frac{L^3}{4\pi^2}
   \left(\frac{2m}{\hbar^2}\right)^{3/2}\epsilon^{1/2}.
 \label{eq:D3}
\end{align}

\section*{[1.3] 三次元箱の臨界温度}
各$\epsilon>0$のBose占有数は$\mu$を0へ近づけるほど増える。凝縮直前は $\mu\to0^-$ であり、連続近似で励起状態の最大粒子数は
\begin{align}
 N_{\rm ex}^{\max}(T)
 &=\int_0^\infty\frac{D_3(\epsilon)}{e^{\epsilon/T}-1}d\epsilon\\
 &=\frac{L^3}{4\pi^2}\left(\frac{2m}{\hbar^2}\right)^{3/2}
   T^{3/2}\int_0^\infty\frac{x^{1/2}}{e^x-1}dx\\
 &=L^3\left(\frac{mT}{2\pi\hbar^2}\right)^{3/2}\zeta\!\left(\frac32\right).
\end{align}
これが $N$ に達する温度が臨界温度である。したがって
\begin{equation}
 \boxed{T_c=\frac{2\pi\hbar^2}{m}
       \left[\frac{N}{L^3\zeta(3/2)}\right]^{2/3}}.
\end{equation}
積分は $\epsilon\to0$ で $\epsilon^{-1/2}$ に比例して可積分である。

\section*{[1.4] 三次元箱の凝縮粒子数}
離散的な基底状態の占有数は $N_0=(e^{-\mu/T}-1)^{-1}$ である。$N_0$ が巨視的になると $\mu\to0^-$ となり、連続近似で数える励起状態は $\mu=0$ での最大数を占める。従って $N=N_0+N_{\rm ex}^{\max}(T)$ であり、
$T<T_c$ では励起粒子数は $N_{\rm ex}(T)=N(T/T_c)^{3/2}$ であり、残りが
基底状態に入るので
\begin{equation}
 \boxed{N_0=N\left[1-\left(\frac{T}{T_c}\right)^{3/2}\right]}.
\end{equation}

\section*{[1.5] 二次元一様系}
二次元箱では運動量空間の円環 $2\pi p\,dp$ と $dp/d\epsilon=m/p$ の積により、
状態密度はエネルギーに依らない定数
$D_2=mL^2/(2\pi\hbar^2)$ となる。$\mu<0$ に対して
\begin{align}
 N_{\rm ex}
 &=D_2\int_0^\infty\frac{d\epsilon}
                 {e^{(\epsilon-\mu)/T}-1}\\
 &=-D_2T\log(1-e^{\mu/T}).
\end{align}
$\mu\to0^-$ でこの数は対数発散する。従って有限温度で全粒子を収容する
最大励起粒子数が有限にならず、この理想一様二次元系では有限温度の凝縮は起こらない。

\section*{[2.1] 二次元調和トラップの状態密度}
エネルギー $\epsilon=n\hbar\omega$ の縮退度は $n+1$ である。エネルギーが
$\epsilon$ 以下の状態数は、大きな $n=\epsilon/(\hbar\omega)$ について
\begin{equation}
 N_{\rm st}(\epsilon)=\sum_{j=0}^{n}(j+1)
 =\frac{(n+1)(n+2)}2
 =\frac{\epsilon^2}{2(\hbar\omega)^2}
   +O\!\left(\frac{\epsilon}{\hbar\omega}\right).
\end{equation}
そのエネルギー微分の最高次項は
\begin{equation}
 \boxed{D_{\rm tr}(\epsilon)=\frac{\epsilon}{(\hbar\omega)^2}}.
\end{equation}

\section*{[2.2] トラップ中の臨界温度}
臨界点では化学ポテンシャルは基底エネルギー $0$ に達する。$\hbar\omega\ll T_c$の連続近似で基底状態を別に扱い、励起状態数の主要項を積分すると、全粒子数は
\begin{align}
 N&=\int_0^\infty\frac{D_{\rm tr}(\epsilon)}{e^{\epsilon/T_c}-1}d\epsilon\\
 &=\left(\frac{T_c}{\hbar\omega}\right)^2
   \int_0^\infty\frac{x}{e^x-1}dx
 =\frac{\pi^2}{6}\left(\frac{T_c}{\hbar\omega}\right)^2.
\end{align}
よって
\begin{equation}
 \boxed{T_c=\frac{\sqrt6}{\pi}\hbar\omega\sqrt N}.
\end{equation}

\newpage
\section*{[2.3] 有限な臨界温度を保つ極限}
与えられたポテンシャルは
$V=V_0(x^2+y^2)/L_0^2$ であり、調和振動子形
$m\omega^2(x^2+y^2)/2$ と係数を比べると
\begin{equation}
 \omega=\frac{1}{L_0}\sqrt{\frac{2V_0}{m}}.
\end{equation}
$V_0$ を固定すると $T_c\propto\sqrt N/L_0$ である。$N\to\infty$ でも
$T_c$ が有限かつ非零に保たれる条件は
\begin{equation}
 \boxed{L_0\propto N^{1/2}}.
\end{equation}

\section*{[2.4] トラップ中の凝縮粒子数}
臨界以下では励起粒子数は $N_{\rm ex}(T)=N(T/T_c)^2$ なので
\begin{equation}
 \boxed{N_0=N\left[1-\left(\frac{T}{T_c}\right)^2\right]}.
\end{equation}
\end{document}
